% GENERATED FILE. Edit canonical lessons, then run npm run generate:books.
% canonical-source-hash: fnv1a64:2db313dcd7eb674a
% canonical-lessons: ZH-C194-malu, ZH-W194-lu-road, ZH-C194-zoulu, ZH-C194-banmaxian, ZH-C194-guangchang, ZH-C194-hongludeng

\chapter{Roads and Crossings}
\label{ch:zh-roads-and-crossings}
\hlchaptermodality{194}

\begin{chapteropening}
\textbf{By the end of this chapter:} \emph{I can talk about roads, going on foot, zebra crossings, squares and traffic lights.}

Complete the last lesson of chapter 194: I can talk about roads, going on foot, zebra crossings, squares and traffic lights.
\end{chapteropening}

\section[\texorpdfstring{\zh{马路} (mǎlù)}{马路 (mǎlù)}]{\zh{马路} (mǎlù) — road; a street for traffic}

\label{lesson:ZH-C194-malu}



\begin{quote}
Before the new one: say the Mandarin for to haggle, then the Mandarin for quality.
\end{quote}

\subsection*{You'll want to know: \zh{马路}}
\textbf{\zh{马路}} — \emph{mǎlù} — \textquotedblleft{}road; a street for traffic\textquotedblright{}.

\begin{quote}
\textbf{\zh{我家在马路对面}} \emph{wǒ jiā zài mǎlù duìmiàn} — My home is across the road.
\end{quote}

The new character, in pieces: \zh{路} — thirteen strokes in two halves: on the left the foot piece, \zh{口} over four short strokes; on the right a short fall, a bar that turns and falls to the left, a long stroke falling to the right, and \zh{口} under them. The next lesson puts it in your hand.

\begin{cousinweb}
\textbf{\zh{马路}} (\emph{mǎlù}) is built from \zh{马} 'horse' and \zh{路} 'road', both inherited from Middle Chinese, \emph{maeX} and \emph{luH} (Guangyun, in Baxter's transcription). Why a road for traffic is named after horses is not settled.
\end{cousinweb}

\begin{grammarlens}[title={What you've built}]
One more word for asking where things are.
\end{grammarlens}

\subsection*{Your turn}
\begin{itemize}
\raggedright
  \item \emph{Say it:} \emph{mǎlù}
  \item \emph{Say it:} \emph{mǎlù}, once more
  \item \emph{Say it:} \emph{wǒ jiā zài mǎlù duìmiàn} — My home is across the road.
  \item \emph{Recall:} say the Mandarin for to haggle, then the Mandarin for quality, then say \emph{mǎlù} again
\end{itemize}

\begin{tcolorbox}[breakable,colback=teal!4,colframe=teal!35!black,title={Before you move on}]
What does \zh{马路} mean? (\textquotedblleft{}Road; a street for traffic\textquotedblright{}.) Say it once more.
\end{tcolorbox}
\section[\texorpdfstring{\zh{路} (lù)}{路 (lù)}]{\zh{路} — the foot piece, then three strokes over \zh{口}}

\label{lesson:ZH-W194-lu-road}



\begin{quote}
The word first, then the character inside it.

\begin{itemize}
\raggedright
  \item \emph{Recall:} say \emph{mǎlù} — \textbf{R1}, one lesson back, before it decays
\end{itemize}

\begin{quote}
\zh{马} \zh{口}
\end{quote}

\emph{Read it:} each one aloud, with nothing said first — \emph{mǎ}, \emph{kǒu}
\end{quote}

\subsection*{Script — the foot piece, then three strokes over \zh{口}}
\hlblockfigure{\detokenize{figures/ZH-W194-lu-road-filmstrip.pdf}}{How \zh{路} is written, stroke by stroke}

\begin{quote}
\zh{路}
\end{quote}

\textbf{Two halves, left before right, and each has \zh{口} in it.}

\textbf{The left half, the foot piece, seven strokes.} \zh{口} at the top, narrow. Under it a short vertical, a short bar from its middle to the right, a short vertical on the left, and a stroke along the foot that rises to the right.

\textbf{The right half, six strokes.} A short stroke falling to the left. A bar that turns and falls down to the left. A long stroke falling down to the right. Then \zh{口} underneath.

7 + 6 = thirteen.

\subsection*{Your turn}
\begin{itemize}
\raggedright
  \item \emph{Write it:} \zh{路}, all thirteen, with the model in front of you
  \item \emph{Cover:} write it again from memory, and count the pieces, not the strokes
\end{itemize}

\begin{tcolorbox}[breakable,colback=teal!4,colframe=teal!35!black,title={Before you move on}]
How many strokes in \zh{路}? (\textbf{Thirteen}.) Where is \zh{口} in the left half? (\textbf{At the top}.) And in the right half? (\textbf{At the bottom}.)

Source: \href{https://raw.githubusercontent.com/chanind/hanzi-writer-data/68d10a4b21150cae5e1ebbd223eed289cf32d90c/data/\%E8\%B7\%AF.json}{Hanzi Writer Data \zh{路}, pinned stroke paths and medians}.
\end{tcolorbox}
\section[zǒulù]{zǒulù — to walk; to go on foot}

\label{lesson:ZH-C194-zoulu}



\begin{quote}
Before the new one: say the Mandarin for quality, then the Mandarin for a street for traffic.
\end{quote}

\subsection*{You'll want to know: zǒulù}
\emph{zǒulù} — \textquotedblleft{}to walk; to go on foot\textquotedblright{}.

\begin{quote}
\emph{wǒ měitiān zǒulù qù shàngbān} — I walk to work every day.
\end{quote}

Its characters come later; for now it is a word for the ear.

\begin{cousinweb}
\emph{zǒulù} is built from \emph{zǒu} and \zh{路} 'road', both inherited from Middle Chinese, \emph{tsuwX} and \emph{luH} (Guangyun, in Baxter's transcription). The dictionary glosses the first 'to hurry'; in Mandarin it has come to mean walking.
\end{cousinweb}

\begin{grammarlens}[title={What you've built}]
One more word for asking where things are.
\end{grammarlens}

\subsection*{Your turn}
\begin{itemize}
\raggedright
  \item \emph{Say it:} \emph{zǒulù}
  \item \emph{Say it:} \emph{zǒulù}, once more
  \item \emph{Say it:} \emph{wǒ měitiān zǒulù qù shàngbān} — I walk to work every day.
  \item \emph{Recall:} say the Mandarin for quality, then the Mandarin for a street for traffic, then say \emph{zǒulù} again
\end{itemize}

\begin{tcolorbox}[breakable,colback=teal!4,colframe=teal!35!black,title={Before you move on}]
What does zǒulù mean? (\textquotedblleft{}To walk; to go on foot\textquotedblright{}.) Say it once more.
\end{tcolorbox}
\section[bānmǎxiàn]{bānmǎxiàn — zebra crossing; pedestrian crossing}

\label{lesson:ZH-C194-banmaxian}



\begin{quote}
Before the new one: say the Mandarin for a street for traffic, then the Mandarin for to go on foot.
\end{quote}

\subsection*{You'll want to know: bānmǎxiàn}
\emph{bānmǎxiàn} — \textquotedblleft{}zebra crossing; pedestrian crossing\textquotedblright{}.

\begin{quote}
\emph{zhèr méiyǒu bānmǎxiàn} — There's no zebra crossing here.
\end{quote}

Its characters come later; for now it is a word for the ear.

\begin{cousinweb}
\emph{bānmǎxiàn} is built from \emph{bānmǎ} 'zebra', literally a mottled horse, and \emph{xiàn} 'a line'. \zh{马} and \emph{xiàn} are inherited from Middle Chinese, \emph{maeX} and \emph{sjenH} (Guangyun, in Baxter's transcription); the dictionary lists \emph{bān} as \emph{paen}, glossed 'mottled' and 'patterned': a crossing striped like a zebra.
\end{cousinweb}

\begin{grammarlens}[title={What you've built}]
One more word for asking where things are.
\end{grammarlens}

\subsection*{Your turn}
\begin{itemize}
\raggedright
  \item \emph{Say it:} \emph{bānmǎxiàn}
  \item \emph{Say it:} \emph{bānmǎxiàn}, once more
  \item \emph{Say it:} \emph{zhèr méiyǒu bānmǎxiàn} — There's no zebra crossing here.
  \item \emph{Recall:} say the Mandarin for a street for traffic, then the Mandarin for to go on foot, then say \emph{bānmǎxiàn} again
\end{itemize}

\begin{tcolorbox}[breakable,colback=teal!4,colframe=teal!35!black,title={Before you move on}]
What does bānmǎxiàn mean? (\textquotedblleft{}Zebra crossing; pedestrian crossing\textquotedblright{}.) Say it once more.
\end{tcolorbox}
\section[guǎngchǎng]{guǎngchǎng — public square}

\label{lesson:ZH-C194-guangchang}



\begin{quote}
Before the new one: say the Mandarin for to go on foot, then the Mandarin for zebra crossing.
\end{quote}

\subsection*{You'll want to know: guǎngchǎng}
\emph{guǎngchǎng} — \textquotedblleft{}public square\textquotedblright{}.

\begin{quote}
\emph{huǒchēzhàn qiánbian yǒu yí ge guǎngchǎng} — There's a square in front of the station.
\end{quote}

Its characters come later; for now it is a word for the ear.

Two third tones meet in this word: guǎng is said with a rising tone before chǎng, and the spelling keeps both marks.

\begin{cousinweb}
\emph{guǎngchǎng} is built from \emph{guǎng} 'wide' and \emph{chǎng} 'an open ground', both inherited from Middle Chinese, \emph{kwangX} and \emph{drjang} (Guangyun, in Baxter's transcription), which glosses the first 'big' and 'broad'.
\end{cousinweb}

\begin{grammarlens}[title={What you've built}]
One more word for asking where things are.
\end{grammarlens}

\subsection*{Your turn}
\begin{itemize}
\raggedright
  \item \emph{Say it:} \emph{guǎngchǎng}
  \item \emph{Say it:} \emph{guǎngchǎng}, once more
  \item \emph{Say it:} \emph{huǒchēzhàn qiánbian yǒu yí ge guǎngchǎng} — There's a square in front of the station.
  \item \emph{Recall:} say the Mandarin for to go on foot, then the Mandarin for zebra crossing, then say \emph{guǎngchǎng} again
\end{itemize}

\begin{tcolorbox}[breakable,colback=teal!4,colframe=teal!35!black,title={Before you move on}]
What does guǎngchǎng mean? (\textquotedblleft{}Public square\textquotedblright{}.) Say it once more.
\end{tcolorbox}
\section[hónglǜdēng]{hónglǜdēng — traffic lights}

\label{lesson:ZH-C194-hongludeng}



\begin{quote}
Before the new one: say the Mandarin for zebra crossing, then the Mandarin for public square.
\end{quote}

\subsection*{You'll want to know: hónglǜdēng}
\emph{hónglǜdēng} — \textquotedblleft{}traffic lights\textquotedblright{}.

\begin{quote}
\emph{lùkǒu yǒu hónglǜdēng hé bānmǎxiàn} — There are traffic lights and a zebra crossing at the junction.
\end{quote}

Its characters come later; for now it is a word for the ear.

\begin{cousinweb}
\emph{hónglǜdēng} is built from \emph{hóng} 'red', \emph{lǜ} 'green' and \emph{dēng} 'a lamp', all inherited from Middle Chinese, \emph{huwng}, \emph{ljowk} and \emph{tong} (Guangyun, in Baxter's transcription).
\end{cousinweb}

\begin{grammarlens}[title={What you've built}]
One more word for asking where things are.
\end{grammarlens}

\subsection*{Your turn}
\begin{itemize}
\raggedright
  \item \emph{Say it:} \emph{hónglǜdēng}
  \item \emph{Say it:} \emph{hónglǜdēng}, once more
  \item \emph{Say it:} \emph{lùkǒu yǒu hónglǜdēng hé bānmǎxiàn} — There are traffic lights and a zebra crossing at the junction.
  \item \emph{Recall:} say the Mandarin for zebra crossing, then the Mandarin for public square, then say \emph{hónglǜdēng} again
\end{itemize}

\begin{tcolorbox}[breakable,colback=teal!4,colframe=teal!35!black,title={Before you move on}]
What does hónglǜdēng mean? (\textquotedblleft{}Traffic lights\textquotedblright{}.) Say it once more.
\end{tcolorbox}
